\chapter{Elliptic fibrations}\label{ch:elliptic}

An elliptic K3 surface is the one kind of K3 surface whose N\'eron--Severi
lattice can be computed by hand from an equation. The tools are Kodaira's
classification of singular fibres, Tate's algorithm for reading the fibre type
off the Weierstrass coefficients, and the Shioda--Tate formula. We state them,
then work the one family this book cares about --- the Inose pencil --- to the
point where the reader can see the lattices $A_2$ and $\lat2\oplus\lat2$ of
Chapter~\ref{ch:lattices} emerge as transcendental lattices.

\section{Weierstrass models}

\begin{definition}
An \emph{elliptic surface} is a smooth projective surface $X$ with a
surjective morphism $\pi\colon X\to C$ to a curve, whose generic fibre is a
smooth curve of genus $1$, with no $(-1)$-curves in fibres (relative
minimality). A \emph{section} is a curve $O\subset X$ with $\pi|_O$ an
isomorphism. If a section exists, the generic fibre is an elliptic curve over
the function field of $C$ and $X$ has a \emph{Weierstrass model}: over
$C=\PP^1$ with coordinate $t$,
\[
  y^2=x^3+a_4(t)\,x+a_6(t),\qquad a_4,a_6\in\C[t],
\]
with discriminant $\Delta=-16(4a_4^3+27a_6^2)$ and $j=1728\cdot 4a_4^3/(4a_4^3+27a_6^2)$.
\end{definition}

\begin{theorem}[\Lstatus]\label{thm:k3-weierstrass}
For $C=\PP^1$, the surface $X$ is K3 if and only if, after making the model
minimal, $a_4$ and $a_6$ have degrees $\le8$ and $\le12$ (as sections of
$\mathcal O(8)$ and $\mathcal O(12)$), not both of lower ``weight'' than
required; then $\Delta$ has degree $24$, and the Euler number
$e(X)=\sum_v e(F_v)=24$ equals the sum of the orders of vanishing of
$\Delta$ at the singular fibres (with $\infty$ included via the degree
count) \cite[Ch.~11]{Huybrechts2016}, \cite[\S~4]{SchuettShioda2010}.
\end{theorem}

\section{Kodaira's classification and Tate's algorithm}

\begin{theorem}[\Lstatus, Kodaira~\cite{Kodaira1963}]\label{thm:kodaira}
A singular fibre of a relatively minimal elliptic surface is one of the types
$I_n$ ($n\ge1$), $II$, $III$, $IV$, $I_n^*$ ($n\ge0$), $IV^*$, $III^*$,
$II^*$ (and, without a section, multiples ${}_mI_n$). The Euler number of
the fibre, the number $m_v$ of its components, and the root lattice spanned
by the components not meeting the zero section are:
\begin{center}
\begin{tabular}{@{}lllll@{}}
\toprule
type & $e(F_v)$ & $m_v$ & root lattice $R_v$ & $\disc R_v$\\
\midrule
$I_n$ & $n$ & $n$ & $A_{n-1}$ & $n$\\
$II$ & $2$ & $1$ & $0$ & $1$\\
$III$ & $3$ & $2$ & $A_1$ & $2$\\
$IV$ & $4$ & $3$ & $A_2$ & $3$\\
$I_n^*$ & $n+6$ & $n+5$ & $D_{n+4}$ & $4$\\
$IV^*$ & $8$ & $7$ & $E_6$ & $3$\\
$III^*$ & $9$ & $8$ & $E_7$ & $2$\\
$II^*$ & $10$ & $9$ & $E_8$ & $1$\\
\bottomrule
\end{tabular}
\end{center}
\end{theorem}

\begin{theorem}[\Lstatus, Tate's algorithm, characteristic $0$~\cite{Tate1975}]\label{thm:tate}
For a minimal Weierstrass model at a place $t=t_0$, the fibre type is read off
the orders of vanishing $(\ord a_4,\ord a_6,\ord\Delta)$ as follows:
\begin{center}
\begin{tabular}{@{}lccc@{}}
\toprule
type & $\ord a_4$ & $\ord a_6$ & $\ord\Delta$\\
\midrule
$I_0$ (smooth) & $\ge0$ & $\ge0$ & $0$\\
$I_n$ & $0$ & $0$ & $n$\\
$II$ & $\ge1$ & $1$ & $2$\\
$III$ & $1$ & $\ge2$ & $3$\\
$IV$ & $\ge2$ & $2$ & $4$\\
$I_0^*$ & $\ge2$ & $\ge3$ & $6$\\
$I_n^*$ & $2$ & $3$ & $n+6$\\
$IV^*$ & $\ge3$ & $4$ & $8$\\
$III^*$ & $3$ & $\ge5$ & $9$\\
$II^*$ & $\ge4$ & $5$ & $10$\\
\bottomrule
\end{tabular}
\end{center}
If $\ord a_4\ge4$ and $\ord a_6\ge6$ the model is not minimal at $t_0$.
\end{theorem}

\begin{warning}
The table is the one place in this book where a small memory slip has large
consequences. An order-$2$ vanishing of $\Delta$ with $a_4,a_6$ nonvanishing
is $I_2$ (an $A_1$), not $I_1$; an order-$4$ vanishing with $a_4\equiv0$ and
$\ord a_6=2$ is $IV$ ($A_2$). Both distinctions decide Picard numbers below.
When in doubt, compute the type from the singularity of the Weierstrass
surface, not from memory --- or check the Euler number sum.
\end{warning}

\section{The Shioda--Tate formula}

\begin{theorem}[\Lstatus, Shioda--Tate~\cite{Shioda1990}]\label{thm:shioda-tate}
Let $X\to\PP^1$ be an elliptic surface with a section $O$ and generic fibre
$F$. Then
\[
  \NS(X)\supseteq \U_{O,F}\oplus\bigoplus_v R_v(-1),\qquad
  \rho(X)=2+\sum_v (m_v-1)+\rk\MW(X),
\]
where $\U_{O,F}$ is the hyperbolic plane spanned by $O$ and $F$ (after the
change of basis $O+F$, $F$), $R_v$ are the fibre root lattices of
Theorem~\ref{thm:kodaira}, and $\MW(X)$ is the Mordell--Weil group of sections.
The sublattice $\U\oplus\bigoplus R_v(-1)$ is the \emph{trivial lattice}; it
has finite index in $\NS(X)$ when $\MW$ is finite, and
$|\disc\NS(X)|=|\disc(\text{trivial})|/|\MW_{\mathrm{tors}}|^2$ in that case.
\end{theorem}

Since $|\disc\NS(X)|=|\disc T(X)|$ for a K3 surface (unimodularity of
$\Lambda$), computing the trivial lattice and the torsion of $\MW$ computes
$\disc T(X)$ --- and for $\rho=20$, that together with positivity pins $T(X)$
down to a short list of binary forms.

\section{The Inose pencil}\label{sec:inose}

\begin{definition}[Inose~\cite{Inose1978}]
For $(\alpha,\beta)\in\C^2$ the \emph{Inose surface} $X_{\alpha,\beta}$ is the
K3 surface with Weierstrass model
\[
  y^2 \;=\; x^3-3\alpha\,t^4\,x+t^5\,(t^2-2\beta t+1).
\]
\end{definition}

\begin{theorem}[\Lstatus, Inose; Shioda~\cite{Shioda2006}]\label{thm:inose}
For elliptic curves $E_1,E_2$ with normalized $j$-invariants $J_i=j(E_i)/1728$,
choose $\alpha,\beta$ with $\alpha^3=J_1J_2$ and $\beta^2=(1-J_1)(1-J_2)$.
Then $X_{\alpha,\beta}$ carries a Shioda--Inose structure with
$T(X_{\alpha,\beta})\cong T(E_1\times E_2)$; in particular
$X_{\alpha,\beta}$ is a double cover of $\Km(E_1\times E_2)$ up to birational
equivalence, and conversely every Kummer surface of a product arises this way.
\end{theorem}

The pencil has fibres of type $II^*$ at $t=0$ and $t=\infty$ for all
$(\alpha,\beta)$: indeed $\ord_0 a_6=5$, $\ord_0 a_4\ge4$, $\ord_0\Delta=10$,
and at $\infty$ (weights $8,12,24$) the same. That accounts for $20$ of the
$24$ in the Euler number and for $E_8\oplus E_8$ in the trivial lattice; the
remaining four are distributed among the roots of $t^2-2\beta t+1$ and the
zeros of $4a_4^3+27a_6^2$ elsewhere.

\begin{example}[\Estatus; the three special members]\label{ex:inose-points}
An exact computation of orders of vanishing (a polynomial computation over
$\Q(i)$, rerunnable in any computer algebra system) gives:
\begin{center}
\begin{tabular}{@{}lll@{}}
\toprule
$(\alpha,\beta)$ & finite places $t_0\neq0$: $(\ord a_4,\ord a_6,\ord\Delta)$ & Euler sum\\
\midrule
$(0,1)$ & $t=1$: $(\infty,2,4)$ & $10+10+4=24$\\
$(1,0)$ & $t=\pm1$: $(0,0,2)$ each & $10+10+2+2=24$\\
$(0,0)$ & $t=\pm i$: $(\infty,1,2)$ each & $10+10+2+2=24$\\
\bottomrule
\end{tabular}
\end{center}
Reading Theorem~\ref{thm:tate} \Lstatus: at $(0,1)$ the fibre at $t=1$ is
$IV$ ($A_2$), so the trivial lattice is $\U\oplus E_8(-1)^2\oplus A_2(-1)$ of
rank $20$ and discriminant $3$; at $(1,0)$ the fibres at $t=\pm1$ are $I_2$
($A_1$ each), trivial lattice of rank $20$ and discriminant $4$; at $(0,0)$ the
fibres are $II$ (no root), trivial lattice of rank $18$. Since $\rho\le20$,
the first two are singular K3 surfaces with $|\disc T|=3$ and $4$ (if
$\MW$ is torsion-free of rank $0$, which one checks separately), hence by
Theorem~\ref{thm:shioda-inose} and the classification of forms,
$T(X_{0,1})\cong A_2$ and $T(X_{1,0})\cong\lat2\oplus\lat2$. The third member
has $\rho\ge18$ from fibres alone; Theorem~\ref{thm:inose} identifies it
($J_1J_2=0$, $(1-J_1)(1-J_2)=0$, so $\{J_1,J_2\}=\{0,1\}$) as the
Shioda--Inose partner of $E_\omega\times E_i$, two non-isogenous CM curves,
for which $\rho=18$ exactly.
\end{example}

The first two are, in Shioda--Inose's notation, the surfaces $X_3$ and $X_4$:
$X_{0,1}$ corresponds to $J_1=J_2=0$, i.e.\ $E_\omega\times E_\omega$ with
$\omega=e^{2\pi i/3}$, and $X_{1,0}$ to $J_1=J_2=1$, i.e.\ $E_i\times E_i$. We
shall meet them again in Chapter~\ref{ch:singular} as the two points of
smallest discriminant on two different modular curves.

\begin{remark}[a recorded slip]
In the program from which this book grew, the point $(0,0)$ was at one stage
labelled ``$X_4$'' and the fibres at $(1,0)$ were at one stage labelled
$I_1$; both errors were caught by the Euler-number and rank counts above and
are recorded in the program's briefs. The table is here so the reader does
not repeat them.
\end{remark}

\section{Exercises}

\begin{exercise}
Verify the entries of Theorem~\ref{thm:kodaira} for $I_n$ and $II^*$ from
the dual graphs (a cycle of $n$ rational curves; the affine $\tilde E_8$
diagram).
\end{exercise}

\begin{exercise}[$\star$]
Compute the orders of vanishing in Example~\ref{ex:inose-points} by hand for
$(\alpha,\beta)=(0,1)$: $a_4=0$, $a_6=t^5(t-1)^2$, $\Delta=-16\cdot27\,t^{10}(t-1)^4$.
\end{exercise}

\begin{exercise}
For $(\alpha,\beta)=(1,0)$, show that $4a_4^3+27a_6^2=27t^{10}\bigl((t^2+1)^2-4t^2\bigr)$
and factor it; confirm $\ord_{t=\pm1}=2$ and that $a_4(\pm1)\neq0$.
\end{exercise}

\begin{exercise}
Using Theorem~\ref{thm:shioda-tate}, show that an elliptic K3 surface with two
$II^*$ fibres has $\rho\ge18$, and that $\rho=20$ requires either two more
fibre components in total or Mordell--Weil rank $2$.
\end{exercise}

\section*{Chapter note}

Theorems~\ref{thm:kodaira}, \ref{thm:tate}, \ref{thm:shioda-tate} are
classical; the tables were checked against Sch\"utt--Shioda's
survey~\cite{SchuettShioda2010} and Silverman's account of Tate's
algorithm~\cite[IV.9]{Silverman1994}. Theorem~\ref{thm:inose} is Inose's with
Shioda's normalization; the relations $\alpha^3=J_1J_2$, $\beta^2=(1-J_1)(1-J_2)$
are as in~\cite{Shioda2006} and should be checked there before being relied
on. The orders of vanishing in Example~\ref{ex:inose-points} are an exact
computation performed for this book; the fibre labels are the table's, and the
identification of $(0,1)$ and $(1,0)$ with $X_3,X_4$ follows from
Theorem~\ref{thm:inose}. An independent computation in the program
(Kuwata--Shioda's model of the same pencil) gave the same orders
$\{10,10,4\}$ at $J=0$ and $\{10,10,2,2\}$ at $J=1$.
